\documentclass[10pt,a4paper]{amsart}
\usepackage[margin=19mm]{geometry}
\usepackage{amsmath,amssymb,mathtools}
\usepackage[scr=rsfs,cal=euler]{mathalfa}
\usepackage{xfrac}
\usepackage[unicode,hidelinks,bookmarks=false]{hyperref}
\newcommand{\ie}{i.e.\ }
\newcommand{\cf}{cf.\ }
\newcommand{\resp}{respectively\ }
\newcommand{\Subsection}{\subsection}
\providecommand{\nobreakdash}{\nobreak\hbox{-}}
\setlength{\emergencystretch}{2em}
\multlinegap0pt
\usepackage{amssymb}
\usepackage{mathtools}
\usepackage{bbm}
\usepackage{xcolor}
\newtheorem{theorem}{Theorem}
\title{Boundedness properties of modified averaging operators and geometrically doubling metric spaces}
\author{J. M. Aldaz, A. Caldera}
\date{Source: arXiv:2403.10445v2 (CC BY 4.0)}
\begin{document}
\maketitle

\noindent\emph{Source and licence.} Boundedness properties of modified averaging operators and geometrically doubling metric spaces. Version arXiv:2403.10445v2 (CC BY 4.0), distributed by arXiv under the Creative Commons Attribution 4.0 International licence, http://creativecommons.org/licenses/by/4.0/, as recorded in the arXiv licence metadata for this exact version and shown on the abstract page https://arxiv.org/abs/2403.10445v2. Full licence notice: \texttt{licenses/arxiv-2403.10445-CC-BY-4.0.txt}.\par\smallskip

\noindent\textit{Editorial scope.} Counted unit: the article's theorem Thrmconj (source0.tex lines 262--264). The article's complete original proof of it is delivered with it (source0.tex lines 266--301, one \texttt{proof} environment of 36 lines). No mathematics is written, repaired or re-derived: the statement and the proof are the article's own lines, in the article's own notation, and no retained line is substituted or re-flowed.  No further source-local context is needed: every cross-reference of every retained line resolves inside the two delivered spans.  Nothing was omitted from any retained span, so the package carries no omission record.  No retained line contains a citation, so the wrapper carries no bibliography and no bibliography entry.  No retained line contains a figure environment, an image-inclusion command or drawing code, so no image is delivered and the package declares no asset dependency.  Editorial formatting only, in addition: an AMS \texttt{amsart} preamble that restates the article's own declarations and macros used by the retained lines, namely the environments \texttt{theorem} on separate counters, together with the standard packages those lines need.  The retained environments print in this document as Theorem 1 in order of appearance, whereas the article prints the same items with the numbering of its own full document.  The complete native source of the article as distributed in the arXiv e-print for this exact version is bundled unchanged as \texttt{sources/156170/source0.tex}.
\begin{theorem}\label{Thrmconj} 
	Let $t > 0$ and let $(X,d,\mu)$ be a metric measure space. The averaging operator $A_{t,r}$ is bounded on $L^1(\mu)$ if and only if $a_{t,r} \in L_{\infty}(\mu)$, in which case $\left \|  A_{t,r}\right \|_{L^1(\mu) \rightarrow L^1(\mu)} = \left \| a_{t,r} \right \|_{\infty}$.
\end{theorem}

\begin{proof}
	Let $0 \le f \in L^1(\mu)$, and suppose  $a_{t,r} \in L^\infty (\mu)$. Since by
	Fubini-Tonelli
	\begin{equation*}\label{fubini}
		\|A_{t,r} f\|_{L^1} 
		=
		\int_X A_{t,r} f(x)  \  d\mu(x)
		=
		\int_X\int_X  \frac{\mathbf{1}_{B(x,rt)}(y) }{\mu (B(x,r))} f (y) \  d\mu(y) \  d\mu(x)
	\end{equation*}
	\begin{equation*}\label{fubini2}
		= 
		\int_X  f (y) \int_X  \frac{\mathbf{1}_{B(y,rt)}(x)}{\mu (B(x,r))}  \  d\mu(x) \  d\mu(y)
		= 
		\int_X  f (y) \ a_{t,r}(y) \  d\mu(y),
	\end{equation*}\\
	it follows from H\"older's inequality that $\|A_{t,r} \|_{L^1(\mu)\to L^1(\mu)} \le \|a_{t,r}\|_\infty. 
	$
	
	
	On the other hand, we claim that if $\|A_{t,r} \|_{L^1(\mu)\to L^1(\mu)} \le C$, then $\|a_{t,r}\|_\infty
	\le C. 
	$
	Towards a contradiction, suppose $C < \|a_{t,r}\|_\infty$ (including the case $\|a_{t,r}\|_\infty
	= \infty$). Then there is a $\lambda > C$ and a measurable set $A_\lambda$ such that 
	$A_\lambda \subset \{ a_{t,r}> \lambda\}$ and $0 < \mu (A_\lambda)  < \infty$. Let $f:= \mathbf{1}_{A_\lambda} \in L^1(\mu)$.
	Then 
	\begin{equation*}\label{fubini3}
		\|A_{t,r} f\|_{L^1} 
		=
		\int_X  f (y) \  a_{t,r}(y) \  d\mu(y)
		>  \int_{A_\lambda}  \lambda  \  d\mu(y)
		= \lambda \mu (A_\lambda) > C \|f\|_{L^1},
	\end{equation*}
thereby contradicting the assumption $\|A_{t,r} \|_{L^1(\mu)\to L^1(\mu)} \le C$.
\end{proof}

\end{document}
